\section{Conclusion and Future Work}

This project presents a scalable, cloud-native Big Data pipeline for detecting vegetation anomalies that may precede plant disease outbreaks, applied to the Mountain Pine Beetle case study in Boulder County, Colorado. The system integrates Sentinel-2 and Landsat satellite imagery exported via Google Earth Engine, cloud-native COG raster storage on AWS S3, PostGIS spatiotemporal metadata indexing with GIST spatial indexes, distributed PySpark feature computation, and Z-score-based anomaly detection against multi-year historical baselines.

Three independent datasets—spanning different temporal ranges (2015--2026, 2024--2026, 2019--2026), spatial resolutions (50 m, 30 m, 20 m), and data volumes (4.3 GB, 2.8 GB, 10.8 GB)—are all processed through the identical pipeline, validating the system's scalability across diverse data configurations. Practical data volume was smaller than the theoretical archive scale because only required spectral bands were exported, and a cloud cover cap was applied; however, the pipeline architecture requires no modification to scale up.

Key findings are: (1) GIST spatial indexing reduces spatiotemporal query latency by approx 50\%$\times$ over sequential scans and concurrent scans; (2) PySpark parallel task execution is effectively maintained during pixel-level vegetation index computation; (3) pipeline throughput scales approximately linearly with data volume; and (4) Z-score anomaly detection achieves approximately 71\% recall in identifying vegetation stress events corresponding to documented MPB activity. Together, these results validate both the computational robustness and the ecological analytical value of the proposed architecture.

Future directions will be to able to compare the results with the ac